\section{Conclusioni}
\label{sec:conclusioni}

L'elaborato si era proposto di confrontare OpenMP, MPI e CUDA sullo stesso
problema, misurandone il comportamento con metriche consolidate e mettendo in
relazione i risultati con i modelli analitici del calcolo parallelo. La tesi
formulata nell'introduzione era che un'unica causa, la variabilità del numero di
iterazioni per pixel, si manifestasse in forme diverse a seconda
dell'architettura: come sbilanciamento del carico fra thread e processi sulla
CPU, come divergenza dei warp sulla GPU. Le sezioni che seguono riassumono che
cosa le misure hanno confermato, che cosa hanno smentito e che cosa è emerso
sul problema.

\subsection{Il problema: un'irregolarità misurabile prima di parallelizzare}

La baseline seriale ha mostrato che il costo per iterazione è costante entro
l'$1{,}3\%$ su tutte le configurazioni misurate: il tempo seriale è proporzionale
al numero totale di iterazioni $\mathcal{W}$, che si può quindi usare al posto del
tempo per ragionare sul lavoro. Il profilo di carico per riga è fortemente
irregolare ($\lambda \approx 2{,}76$, $\mathrm{CoV} \approx 0{,}97$), e questa
irregolarità dipende dalla geometria dell'insieme più che dalla soglia iterativa:
portando $N_{\max}$ da $500$ a $5000$, $\lambda$ cresce del $2{,}9\%$, mentre il
lavoro cresce di un fattore $9{,}57$ e solo lo $0{,}116\%$ dei pixel cambia
classificazione.

La conseguenza metodologica è il risultato più generale del lavoro: il profilo
di carico e la matrice degli escape time, ottenuti da una sola esecuzione
seriale, bastano a formulare previsioni quantitative sulle versioni parallele
prima di scriverle. Sulla CPU queste previsioni si sono rivelate accurate; sulla
GPU descrivono correttamente il lavoro, ma non il tempo.

\subsection{CPU: la regola di assegnamento decide lo speedup}

Sulla CPU le previsioni ricavate dal solo profilo seriale sono state confermate
(sezione~\ref{sec:validazione-previsioni}): $19$ dei $25$ punti di misura
rientrano nell'$1\%$ del limite predetto. I sei restanti, tutti a $p \ge 16$, se
ne discostano fra il $2{,}6$ e l'$8{,}5\%$, per costi che il modello non include.
La regola a blocchi, realizzata con \texttt{schedule(static)}, raggiunge il
proprio limite entro lo $0{,}6\%$: la perdita del $62\%$ di efficienza a $32$
thread è dovuta quasi per intero alla contiguità delle righe costose, che
finiscono nello stesso blocco. Cambiando soltanto la regola, senza aggiungere
comunicazione, la decomposizione ciclica resta entro l'$1\%$ dall'ideale fino a
$8$ processi; a $16$ e $32$ se ne allontana, per costi seriali su rank $0$ ---
il riordino di de-interlacciamento e un secondo buffer da $4$~MiB --- che la sua
$e(p)$ costante quantifica in circa $1{,}4$~ms, e a $p=16$ per una gather anomala
($1{,}66$~ms contro i $0{,}36$~ms del block).
Le regole dinamiche a grana fine arrivano a $S = 30{,}99$ (OpenMP
\texttt{dynamic,1}) e $S = 31{,}10$ (MPI master--worker) su $32$ unità. Nello
sweep MPI il rapporto fra lo speedup peggiore (block) e il migliore
(master--worker) è $2{,}66$, pari allo sbilanciamento $\lambda_{32}$ previsto per
i blocchi.

Il dato che accomuna i due paradigmi è che la previsione dipende dalla
\emph{regola} e non dal meccanismo che la realizza: una clausola
\texttt{schedule(static)} in memoria condivisa e una decomposizione esplicita con
gather in memoria distribuita cadono sullo stesso limite entro l'$1\%$ fino a
$p = 16$. Anche il confronto fra i due scheduler dinamici lo conferma: fino a
$p = 16$ il passaggio da memoria condivisa a scambio di messaggi costa meno del
$2\%$.

\subsection{GPU: la stessa irregolarità pesa poco sul tempo}

Una sola L40S calcola l'immagine di riferimento $128$ volte più velocemente di un
core e $4{,}11$ volte più velocemente del miglior risultato CPU, il master--worker
MPI su $32$ unità. Tutte le geometrie del blocco riproducono il checksum della
baseline, e la geometria stessa si è rivelata una leva di second'ordine: fra la
più veloce e la più lenta corre l'$8{,}2\%$.

Le grandezze calcolabili dalla matrice degli escape time sono esatte, ma spiegano
poco delle differenze di tempo. L'occupancy vale $1$ per tutte le geometrie,
perché il kernel non esaurisce i registri. La perdita di coalescenza non comporta
una penalità misurabile: la $1{\times}256$, che la subisce, è la geometria più
veloce, e il traffico di memoria resta sotto il $3\%$ della banda anche nel caso
peggiore. Una riduzione prevista del $7{,}2\%$ dei cicli eseguiti si traduce in
un guadagno misurato dello $0{,}4\%$. L'unica differenza riconducibile alla
distribuzione del carico è quella fra le geometrie orizzontali pure, circa il
$5\%$, interpretata come effetto della granularità con cui il lavoro si
ripartisce fra gli SM. L'irregolarità del carico, che sulla CPU spiega quasi per
intero le prestazioni, sulla GPU ne spiega quindi soltanto pochi punti
percentuali.

Lo Sweep~B aggiunge tre elementi. La scheda si satura da circa $2048^2$ pixel,
dove il throughput si stabilizza attorno a $440$~GFLOP/s e lo speedup rispetto a
un core a $152$. Alle taglie inferiori il rendimento cala, e né il trasferimento
né un costo fisso di lancio bastano a spiegarlo: il calo è compatibile con il
sottoutilizzo degli SM, che senza performance counter non è stato possibile
verificare. La copia dei risultati pesa fra lo $0{,}8$ e il $7{,}1\%$ del tempo e
si ammortizza al crescere di $N_{\max}$. Infine, la posizione del processo
rispetto alla GPU non è garantita: nel job misurato la copia attraversava i due
socket e raggiungeva una banda asintotica di circa $24$~GB/s, tre quarti di
quella nominale del PCIe, senza che fosse possibile isolare il contributo del
passaggio fra i socket.

\subsection{Le ipotesi iniziali alla prova dei dati}

\paragraph{La tesi dell'unica causa.}
È confermata come principio che organizza l'intero lavoro, ma con un peso molto
diverso nei due casi. Sulla CPU la variabilità degli escape time determina lo
speedup, e lo fa in modo prevedibile: le regole statiche saturano il limite
predetto entro lo $0{,}6\%$, e gli scostamenti maggiori corrispondono a costi
esterni al modello. Sulla GPU la stessa variabilità è misurabile come divergenza
dei warp, ma il suo effetto sul tempo è marginale nell'intervallo esplorato.
Anche la conclusione più suggestiva, che la stessa causa richieda rimedi opposti
--- mescolare il carico sulla CPU, raggruppare pixel simili sulla GPU --- resta
per la GPU una previsione del modello: il raggruppamento sparso non è
realizzabile con una forma di blocco e non è stato misurato, e lo Sweep~A indica
che il tempo è poco sensibile ai cicli dei warp.

\paragraph{Amdahl e Karp--Flatt.}
La legge di Amdahl, con una frazione seriale trascurabile, prevede uno speedup
quasi ideale e da sola non spiega le perdite osservate; il valore di $1-P$ resta
peraltro stimato e non misurato, e sono i tempi paralleli a vincolarlo
dall'alto (sezione~\ref{sec:parallelizzazione}). La metrica di Karp--Flatt
segnala correttamente le perdite, ma quasi mai misura una frazione seriale: per
la regola a blocchi $e(p)$ decresce con $p$ e segue la forma
$(\lambda_p-1)/(p-1)$, cioè lo sbilanciamento riscritto in un'altra variabile. È
diagnostico il suo andamento, più che il suo valore: una $e(p)$ decrescente
indica uno sbilanciamento il cui limite cresce con $p$; una crescente o un
overhead che peggiora con il parallelismo, come la contesa sulla coda di
\texttt{dynamic,1}, oppure un tetto fisso di granularità, come in
\texttt{dynamic,16} e \texttt{dynamic,64} --- a distinguerli è se la curva
converga a $1/S_{\max}$ o resti ordini di grandezza al di sotto. L'unica
eccezione, e l'unico caso della campagna in cui la metrica misura davvero codice
seriale, è la decomposizione ciclica MPI, la cui $e(p)$ resta costante attorno a
$0{,}002$.

\paragraph{Andamenti più che valori assoluti.}
L'introduzione aveva anticipato che i valori assoluti, legati alle risorse
disponibili, contassero meno degli andamenti. I risultati lo confermano. Le
previsioni sulla CPU sono rapporti, e non dipendono dal valore di $T_1$. Le sei
baseline seriali dei diversi job concordano entro lo $0{,}8\%$, e due schede
L40S fisicamente diverse hanno dato tempi entro lo $0{,}7\%$: nessuna di queste
differenze modifica le conclusioni.

\subsection{Risultati inattesi}

Alcuni risultati non erano previsti, e sono fra i più istruttivi.
\begin{itemize}
    \item \textbf{La granularità può annullare il bilanciamento dinamico.}
    \texttt{dynamic,64} si ferma a $S = 6{,}17$ già da $8$ thread: con soli $16$
    chunk il tempo non può scendere sotto quello del chunk più pesante, e lo
    scheduling dinamico eredita il limite di una decomposizione statica a $16$
    blocchi. La stessa regola, applicata ai $64$ chunk di \texttt{dynamic,16},
    predice $S_{\max} = 23{,}64$ contro i $23{,}21$ misurati: il tetto di
    granularità, e non il costo della coda, spiega anche quella politica.
    \item \textbf{Una politica adattiva può lavorare contro il problema.}
    \texttt{guided} perde circa il $31\%$ di efficienza a ogni grado di
    parallelismo, perché assegna blocchi contigui grandi proprio all'inizio,
    quando il carico andrebbe compensato.
    \item \textbf{Il modello di memoria pesa poco.} A $32$ unità il master--worker
    MPI è più veloce dello scheduler dinamico di OpenMP dell'$1{,}1\%$; la contesa
    sulla coda condivisa di OpenMP è una spiegazione plausibile, non verificata.
    A parità di core allocati, però, OpenMP torna avanti di circa il $2\%$.
    \item \textbf{Uno speedup superlineare solo apparente.} A $p = 2$ tutte le
    politiche OpenMP superano di poco lo speedup ideale ($S \approx 2{,}011$). La
    causa è la variabilità di $T_1$ fra job distinti: la baseline di un altro job,
    più veloce di $3{,}3$~ms, riporta lo speedup a $2{,}002$. Per questo ogni
    misura è riferita alla baseline del proprio job, misurata sul nodo nello
    stesso stato.
    \item \textbf{Sulla GPU le leve attese spiegano poco.} Occupancy, coalescenza
    e divergenza, individuate in fase di progettazione come fattori determinanti,
    non spiegano le differenze di tempo fra le geometrie. Restano aperti il
    primato della geometria $1{\times}256$ e il ritardo delle forme
    bidimensionali coalescenti, che senza accesso ai performance counter non è
    stato possibile indagare.
\end{itemize}

\subsection{Limiti e sviluppi}

I risultati vanno letti alla luce dell'ambiente di esecuzione
(sezione~\ref{sec:ambiente}). Tutte le misure provengono da un unico nodo: MPI è
stato eseguito su memoria condivisa, e la frazione di comunicazione misurata, mai
superiore al $3{,}4\%$, è un limite inferiore ottimistico rispetto a
un'esecuzione distribuita, nella quale il master--worker potrebbe perdere il
vantaggio sul ciclico. Il nodo è condiviso con altri utenti, e il minimo sulle
ripetizioni ne ha assorbito il rumore; anche la posizione del processo rispetto
alla GPU dipende dall'assegnazione di SLURM e non è controllabile. Sulla GPU tutte
le misure usano una sola scheda. La doppia precisione, imposta dalla
riproducibilità, colloca il kernel su un picco nominale $64$ volte inferiore a
quello FP32; i picchi, inoltre, sono documentati per la L40 e non misurati.
L'impossibilità di accedere ai performance counter impedisce infine di spiegare
i risultati residui.

Restano come sviluppi naturali il Buddhabrot, proposto nell'introduzione per
osservare il regime dominato dalla contesa sulla memoria condivisa, e non
affrontato; le versioni ibride, che combinano MPI con OpenMP o con CUDA; un
esperimento in singola precisione, per quantificare il compromesso fra
precisione e throughput; il cronometraggio separato dell'allocazione della
matrice, che fisserebbe per misura la frazione seriale $1-P$ oggi solo stimata;
la misura diretta del costo di riordino nella
decomposizione ciclica MPI, insieme alla rimisura del punto a $p=16$; l'accesso
ai performance counter, per indagare le due questioni aperte dello Sweep~A; e
un'esecuzione su più nodi, che darebbe alla frazione di comunicazione il suo
valore reale.
